\chapter{Introdução à Probabilidade}

Em probabilidade, estudam-se os experimentos aleatórios. Assim, primeiro, é necessário definir um experimento aleatório, diferenciando o mesmo de um experimento determinístico.

\begin{definition}
	Os \emphasis{Experimentos Determinísticos} são experimentos que ao serem repetidos "nas mesmas condições" conduzem ao mesmo resultado.
\end{definition}

\begin{example}
	Se aquecermos a água a 100 $^{\text{o}}$C, ela entrará em ebulição.
\end{example}

\begin{definition}
	Os \emphasis{Experimentos Aleatórios} são experimentos que ao serem repetidos "nas mesmas condições" em geral, não produzem o mesmo resultado.
\end{definition}

\begin{example}
	São exemplos de experimentos aleatórios:
	\begin{itemize}
		\item O número de grãos produzidos por hectares;
		\item Tipos de cultivares de milho;
		\item Tempo de uma reação química;
		\item Peso de um animal.
	\end{itemize}
\end{example}

\section{Espaço Amostral e Eventos}

\begin{definition}
	Denominaremos espaço amostral associado a um experimento ($E$) o conjunto de seus possíveis resultados, que representaremos por $S$, cujos elementos serão denominados eventos simples ou pontos amostrais.
\end{definition}

\begin{pageWidthArea}
	\begin{example}
		\label{example:pecas_boas_defeituosas}
		\begin{multicols}{2}
			\begin{enumerate}\hfill
				\item Seja o experimento $E$: Lançamento de uma moeda. Então, $S=\{C,\text{ }K\}$, onde $C$ representa cara e $K$ coroa.
			
				\item Seja o experimento $E$: três peças são retiradas de uma linha de produção e classificadas em Boa ($B$) ou Defeituosas ($D$). Assim,
					\label{example:item:pecas_boas_defeituosas}
					\begin{center}
						\(S=\{BBB\), \(BBD\), \(BDB\), \(DBB\), \(BDD\), \(DBD\), \(DBD\), \(DDB\), \(DDD\}\).
					\end{center}
		
				\item Seja o experimento $E$: número de plantas doentes em um tomateiro. Então,
					\[
						S=\{0, 1, 2, \dots, n\}\text{,}
					\]
					onde $n$ é o número máximo de plantas neste tomateiro.
					
				\columnbreak
			
				\item Seja $E$: uma moeda é lançada sucessivamente até que ocorra cara pela primeira vez. Assim,
					\[
						S=\{C, KC, KKC, KKKC, KKKKC, \dots \}\text{.}
					\]
			
				\item Seja $E$: tempo de vida de um animal. Então,
					\[
						S=\{t\in\mathbb{R}/ t\geqslant 0\}\text{,}
					\]
					ou
					\[
						S=\mathbb{R}_{+}\text{.}
					\]
			\end{enumerate}
		\end{multicols}
	\end{example}
\end{pageWidthArea}

\begin{definition}
	Denominaremos de evento a todo resultado ou subconjunto de resultados de um experimento.
\end{definition}

\begin{remark}
	Eventos simples são representados por um conjunto unitário.
\end{remark}

\begin{remark}
	Diremos que um evento $A$ ocorre quando o resultado do experimento é um evento simples pertencente ao evento $A$.
\end{remark}

\begin{example}
	Consideremos o Item \ref{example:item:pecas_boas_defeituosas} do Exemplo \ref{example:pecas_boas_defeituosas} e, seja um evento $A$: "Duas peças são boas". Tem-se
	\[
		A=\{BBD, BDB, DBB\}\text{,}
	\]
	onde $B$: Boa e $D$: Defeituosa. Então, $A$ ocorre se um dos três eventos simples:
	\begin{center}
		$BBD$, $BDB$, $DBB$,
	\end{center}
	ocorrer.
\end{example}

\begin{example}
	Considere o lançamento de duas moedas e, seja o evento B: "o número de caras é igual ao número de coroas". Assim,
	\[
		B=\{CK, KC\}\text{,}
	\]
	onde $C$: cara e $K$: coroa.
\end{example}

\section{Operação entre Eventos}

\begin{definition}
	\label{def:reuniao_dois_eventos}
	A reunião de dois eventos $A$ e $B$, denotada por $A\cup B$, é o evento que ocorre se pelo menos um deles ocorre;
\end{definition}

\begin{definition}
	\label{def:intercecao_dois_eventos}
	A intersecção de dois eventos $A$ e $B$, denotada por $A\cap B$, é o evento que ocorre quando $A$ não ocorre.
\end{definition}

\begin{definition}
	\label{def:complementar_evento}
	O complementar do evento $A$, denotado por $A^{c}$, é o evento que ocorre quando $A$ não ocorre.
\end{definition}

\begin{example}
	Um fazendeiro com o intuito de fazer uma amostragem simples, enumerou de um a quinze, quinze gado Nelore. Sejam $A$ e $B$ os seguintes eventos:
	\begin{itemize}
		\item $A$: "o número do gado sorteado é par" e
		\item $B$: "o número do gado sorteado é múltiplo de 3".
	\end{itemize}
	
	Determine os eventos $A\cup B$, $A\cap B$ e $A^{c}$.\\
	
	Escrevendo o espaço amostral $S=\{1$, $2$, $3$, $4$, $5$, $6$, $7$, $8$, $9$, $10$, $11$, $12$, $13$, $14$, $15\}$ e os eventos $A=\{2$, $4$, $6$, $8$, $10$, $12$, $14\}$ e $B=\{3$, $6$, $9$, $12$, $15\}$, pode-se determinar os conjuntos
	\begin{itemize}
		\item $A\cup B=\{2$, $3$, $4$, $6$, $8$, $9$, $10$, $12$, $14$, $15\}$,
		\item $A\cap B=\{6$, $12\}$ e
		\item $A^{c}=\{1$, $3$, $5$, $7$, $9$, $11$, $13$, $15\}$.
	\end{itemize}
\end{example}

\begin{remark}
	Os eventos $A$ e $B$ são ditos mutuamente exclusivos, se eles não podem ocorrer simultaneamente. Isto é $A\cap B=\emptyset$.
\end{remark}

\begin{theorem}
	A partir das Definições \ref{def:reuniao_dois_eventos}, \ref{def:intercecao_dois_eventos} e \ref{def:complementar_evento}, tem-se as seguintes propriedades:
	\begin{enumerate}
		\item $(A\cup B)\cap C = (A\cap C)\cup (B\cap C)$;
		\item $(A\cap B)\cup C = (A\cup C)\cap (B\cup C)$;
		\item $(A\cup B)^c = A^c \cap B^c$;
		\item $(A\cap B)^c= A^c\cup B^c$.
	\end{enumerate}
	
	\begin{proof}
		Essas propriedades são propriedades de conjuntos e, portanto, não serão demonstradas aqui. \todo{Adicionar alguma referência aqui}.
	\end{proof}
\end{theorem}

\begin{definition}
	O evento $\displaystyle\bigcup_{i=1}^{n} A_i$ é o evento que ocorre quando pelo menos um dos eventos $A_i$ ocorre, $i=1,2,\dots, n$.
\end{definition}

\begin{definition}
	O evento $\displaystyle\bigcap_{i=1}^{n} A_i$ é o evento que ocorre quando todos os eventos $A_i$ ocorrem simultaneamente, $i=1,2,\dots, n$.
\end{definition}

\section{Definição Clássica de Probabilidade}

Seja $S$ com $n$ eventos simples igualmente possíveis e seja $A$ um evento de $S$ formado de $m$ eventos simples. A probabilidade de ocorrer $A$, é definida por
\begin{equation}
	P(A)=\frac{m}{n}\text{.}
\end{equation}

\begin{example}
	Joga-se um dado, qual a probabilidade de ocorrer:
	\begin{enumerate}[label=\alph*)]
		\item Evento A: "Face 4"\\
		
		tem-se $A=\{4\}$, logo $m=1$ e $n=6$, pois $S=\{1$, $2$, $3$, $4$, $5$, $6\}$, então
		\[
			P(A)=\frac{1}{6}\text{.}
		\]
		
		\item Evento B: "Face Ímpar"\\
		
		tem-se $B=\{1$, $3$, $5\}$, logo $m=3$, então
		\[
			P(B)=\frac{3}{6}=\frac{1}{2}\text{.}
		\]
	\end{enumerate}
\end{example}

\section{Definição Frequentista de Probabilidade}

Repete-se o experimento $n$ vezes e anota quantas vezes o evento $A$ associado a esse experimento ocorre. Assim a probabilidade de ocorrer o evento $A$ é dado por:
\begin{equation}
	P(A)=\frac{n(A)}{n}\text{,}
\end{equation}
onde $n(A)$ é o número de vezes que o evento $A$ ocorre.

\begin{sidepicture}{0cm}{table}{AAAA}
	\label{table:resistencia_choque_arranhao}
	\begin{tabular}{ccc}\\\toprule
		\multirow{2}{*}{\begin{tabular}[c]{@{}c@{}}Resistência\\ a arranhão\end{tabular}} & \multicolumn{2}{c}{Resistência a choque} \\ \cline{2-3} 
			  & alta & baixa\\ \midrule
		alta  & 80   & 9 \\
		baixa & 6    & 5 \\ \bottomrule
	\end{tabular}
\end{sidepicture}

\begin{example}
	Amostras de plástico policarbonato são analisadas com relação à resistência a arranhões e a choque. Considerando os dados da Tabela \ref{table:resistencia_choque_arranhao}, calcule a probabilidade de uma unidade ter:
	
	\begin{enumerate}[label=\alph*)]
		\item baixa resistência a arranhão (bra):
		\[
			P(bra)=\frac{n(bra)}{n}=\frac{11}{100}=0,11\text{.}
		\]
		
		\item alta resistência a choque (arc):
		\[
			P(arc)=\frac{n(arc)}{n}=\frac{86}{100}=0,86\text{.}
		\]
		
		\item baixa resistência em ambos:
		\[
			P(bra\cap brc)=\frac{n(bra\cap brc)}{n}=\frac{5}{100}=0,05\text{.}
		\]
	\end{enumerate}
\end{example}

\section{Definição Axiomática de Probabilidade}

A probabilidade será definida numa classe de eventos ($A$) do espaço amostral ($S$) que satisfaz certas propriedades:
\begin{enumerate}
	\item $0\leqslant P(A)\leqslant 1$;
	\item $P(S)=1$;
	\item $P(\emptyset)=0$;
	\item $P(A^c)=1-P(A)$;
	\item $P(A\cup A^c)=P(S)$;
	\item $P\left ( \displaystyle \bigcup_{i=1}^{n} A_i \right ) = \displaystyle\sum_{i=1}^{n}P(A_i)$, com os $A_i$ mutuamente exclusivos.
\end{enumerate}

\section{Regra de adição de probabilidade}

Sejam $A$ e $B$ eventos de $S$, então a probabilidade de que ocorra um destes dois eventos, $A$ ou $B$ (Figura \ref{picture:eventos_intersessao}), é dada por
\begin{equation}
	P(A\cup B)=P(A)+P(B)-P(A\cap B)\text{.}
\end{equation}

\begin{remark}
	Se os dois eventos, $A$ e $B$, forem mutuamente exclusivos (Figura \ref{picture:eventos_mutuamente_exclusivos}), então
	\begin{equation}
		P(A\cup B)=P(A)+P(B)\text{,}
	\end{equation}
	pois $A\cap B=\emptyset$, então, $P(A\cap B)=0$.
\end{remark}

\begin{sidepicture}{5cm}{picture}{QQQQQ.}
	\label{picture:eventos_intersessao}
	\begin{tikzpicture}
		\scope
			\draw[pattern=horizontal lines, distance=25pt, pattern color=ocre!50] (-0.3,0) circle (1);
		\endscope
		% right hand
		\scope
			\draw[pattern=vertical lines, distance=25pt, pattern color=ocre!50] (1.3,0) circle (1);
		\endscope
		% outline
		\draw[color=ocre] (-0.3,0) circle (1) (-0.3,1)  node [text=ocre,above] {$A$}
		      (1.3,0) circle (1) (1.3,1)  node [text=ocre,above] {$B$}
		      (-2,-2) rectangle (3,2) node [text=black,above] { };
	\end{tikzpicture}
\end{sidepicture}

\begin{sidepicture}{0cm}{picture}{QQQQQ.}
	\label{picture:eventos_mutuamente_exclusivos}
	\begin{tikzpicture}
		\scope
			\draw[pattern=horizontal lines, distance=25pt, pattern color=ocre!50] (-0.7,0) circle (1);
		\endscope
		% right hand
		\scope
			\draw[pattern=vertical lines, distance=25pt, pattern color=ocre!50] (1.7,0) circle (1);
		\endscope
		% outline
		\draw[color=ocre] (-0.7,0) circle (1) (-0.7,1)  node [text=ocre,above] {$A$}
		      (1.7,0) circle (1) (1.7,1)  node [text=ocre,above] {$B$}
		      (-2,-2) rectangle (3,2) node [text=black,above] { };
	\end{tikzpicture}
\end{sidepicture}

\begin{example}
	A probabilidade de um indivíduo ser obeso num grupo de pacientes é de 30\% e de ser sedentário é de 60\%. Sabe-se, ainda, que a probabilidade de um individuo ser obeso e sedentário é de 28\%. Determine a probabilidade de selecionar um individuo desse grupo e ele:
	
	\begin{enumerate}[label=\alph*)]
		\item ser obeso ou sedentário.
		
			\begin{align*}
				P(A\cup B)&=P(A)+P(B)-P(A\cap B)\\
						  &=0,3+0,6-0,28\\
						  &=0,9-0,28\\
						  &=0,62\text{.}
			\end{align*}
			
		\item Não ser obeso e nem sedentário.
		
			\begin{align*}
				P(A^c\cap B^c)&=P((A\cup B)^c)\\
							  &=1-P(A\cup B)\\
							  &=1-0,62\\
							  &=0,28\text{.}
			\end{align*}
	\end{enumerate}
\end{example}

\begin{example}
	No atendimento de uma fila são considerados dois casos: $A$: "grávidas" e $B$: "pessoas com mais de 60 anos". Seja $P(A)=0,06$ e $P(B)=0,12$. Qual é a probabilidade de que seja necessário dar prioridade no atendimento?
	
	Considerando $A\cap B=\emptyset$, tem-se $P(A\cap B)=0$, e, então
	\begin{align*}
		P(A\cup B)&=P(A)+P(B)\\
				  &=0,06+0,12\\
				  &=0,18\text{.}
	\end{align*}
\end{example}

\section{Probabilidade Condicional}

Dados dois eventos $A$ e $B$, a probabilidade condicional de $A$ dado que ocorreu $B$ é representada por $P(A\setminus B)$ e é dada por
\begin{equation}
	P(A\setminus B)=\frac{P(A\cap B)}{P(B)}\text{,}
\end{equation}
com $P(B)>0$.

\begin{example}
	Em um lançamento de um dado, de 6 faces, $F_1$, $F_2$, $F_3$, $F_4$, $F_5$ e $F_6$, tem-se a probabilidade de se obter a face 2:
	\[
		P(F_2)=\frac{1}{6}\text{.}
	\]
	
	A probabilidade, entretanto, de se ter a face 2, visto que o lançamento terminou com face par, é dado por
	\begin{align*}
		P(F_2\setminus F_{\text{par}})&=\frac{P(F_2 \cap F_{\text{par}})}{P(F_{\text{par}})}\\
									   &=\frac{P(F_2)}{P(F_{\text{par}})}\\
									   &=\frac{1/6}{1/2}\\
									   &=\frac{1}{6}\cdot \frac{2}{1}\\
									   &=\frac{1}{3}\text{.}
	\end{align*}
\end{example}

\newpage

\begin{pageWidthArea}
	\begin{pageWidthAreaPicture}{frame}{Quadro contendo dados de filmes para o Exemplo.}
		\label{quadro:dados_filmes_exemplo_pc}
		\hspace{-15pt}
		\resizebox{\textwidth}{!}{%
		\begin{tabular}{|c|c|c|c|c|}
			\hline
			Filme     & \multirow{2}{*}{\emphasis{A}ventura} & \multirow{2}{*}{\emphasis{C}omédia} & \multirow{2}{*}{\emphasis{R}omance} & \multirow{2}{*}{Total} \\ \cline{1-1}
			Sexo      &                           &                          &                          &                        \\ \hline
			\emphasis{M}asculino & 50                        & 30                       & 10                       & 90                     \\ \hline
			\emphasis{F}eminino  & 20                        & 30                       & 60                       & 110                    \\ \hline
			Total     & 70                        & 60                       & 70                       & 200                    \\ \hline
		\end{tabular}%
		}
	\end{pageWidthAreaPicture}
	
	\vspace{20pt}

	\begin{example}
		\begin{multicols}{2}
			Analisando os dados do Quadro \ref{quadro:dados_filmes_exemplo_pc}, é possível determinar\\
	
			\begin{enumerate}[label=\alph*), itemsep=10pt]
				\item $P(A)=\dfrac{70}{200}=0,35$.
		
				\item $P(A\setminus M)=\dfrac{50}{90}=0,5556$, ou, ainda,
				\begin{align*}
					P(A\setminus M)&=\frac{P(A\cup M)}{P(M)}\\
							&=\frac{50}{200}\cdot \frac{200}{90}\\
							&=\frac{50}{90}\\
							&=0,5556\text{.}
				\end{align*}
		
				\item $P(M\setminus A)=\dfrac{50}{70}=0,7143$.
		
				\item $P(R^c)=1-P(R)=1-\dfrac{70}{200}=0,65$.
		
				\item $P(F\cap C)=\dfrac{30}{200}=0,15$.
				
				\item $P(F\cup A) = P(F) + P(A) - P(F\cap A)$, assim,
				\begin{align*}
					P(F\cup A) &= \frac{110}{200}+\frac{70}{200}-\frac{20}{200}\\
							   &= \frac{110+70-20}{200}\\
							   &= \frac{160}{200}\\
							   &=0,80\text{.}
				\end{align*}
				
				\item $P(M\setminus R^c)$, que é dado por
				\begin{align*}
					P(M\setminus R^c) &= \frac{P(M\cap R^c)}{P(R^c)}\\
									  &= \frac{80}{200}\cdot \frac{200}{130}\\
									  &= \frac{80}{130}\\
									  &= 0,6154\text{.}
				\end{align*}
			\end{enumerate}
		\end{multicols}
	\end{example}
\end{pageWidthArea}

\section{Regra do Produto}

Sabendo que
\[
	P(A\setminus B)=\frac{P(A \cap B)}{P(B)}\text{,}
\]
tem-se
\begin{equation}
	P(A\cap B)=P(A\setminus B)\cdot P(B)\text{.}
\end{equation}

\begin{example}
	A probabilidade de que uma bateria de automóvel, sujeita a alta temperatura no compartimento do motor, sofra baixa corrente de carga é de 0,7. A probabilidade da bateria estar sujeita a alta temperatura no comportamento do motor é de 0,05. Determine a probabilidade da bateria estar sujeita a baixa corrente de carga e a alta temperatura do motor.
	
	Sejam os eventos BC: "baixa corrente"; AT: "Alta temperatura", então $P(BC\cap AT)=P(BC\setminus AT)\cdot P(AT)$, ou seja,
	\begin{align*}
		P(BC\cap AT)&=0,7\cdot 0,05\\
					&=0,035\text{.}
	\end{align*}
\end{example}

\section{Independência de Eventos}

Dois eventos $A$ e $B$ são independentes se a informação da ocorrência ou não de $B$ não altera a probabilidade da ocorrência de $A$, isto é:
\[
	P(A\setminus B)=P(A)\text{,}
\]
ou, ainda,
\[
	P(A\cap B)=P(A)\cdot P(B)\text{.}
\]

\begin{sidepicture}{6cm}{picture}{AAAA}
	\label{picture:prob_independ_ex1}
	\begin{tikzpicture}[
		grow=right,
		level 1/.style ={level distance=2cm, sibling distance=2.6cm},
		level 2/.style ={level distance=2cm, sibling distance=1.2cm},
	]
		\node [prob_tree_origin] {}
		child {
			node [prob_tree_node] {$V_1$}
			child {
				node [prob_tree_node] {$V_2$}
				edge from parent
				node[below, prob_tree_text]  {$\frac{2}{4}$}
			} child {
				node [prob_tree_node] {$B_2$}
				edge from parent
				node[above, prob_tree_text]  {$\frac{2}{4}$}
			}
			edge from parent
			node[below, prob_tree_text]  {$\frac{3}{5}$}
		} child {
			node [prob_tree_node] {$B_1$}
			child {
				node [prob_tree_node]{$V_2$}
				edge from parent
				node[below, prob_tree_text]  {$\frac{3}{4}$}
			} child {
				node [prob_tree_node]{$B_2$}
				edge from parent
				node[above, prob_tree_text]  {$\frac{1}{4}$}
			}
			edge from parent
			node[above, prob_tree_text]  {$\frac{2}{5}$}
		};
	\end{tikzpicture}
\end{sidepicture}

\begin{sidepicture}{0cm}{table}{AAAA}
	\label{table:prob_independ_ex1}
	\begin{tabular}{cc}\\\toprule
		Resultados & Probabilidades\\  \midrule
		$B_1 B_2$ & \vspace{5pt}$\dfrac{2}{5}\cdot\dfrac{1}{4}=\dfrac{2}{20}$ \\ 
		$B_1 V_2$ & \vspace{5pt}$\dfrac{2}{5}\cdot\dfrac{3}{4}=\dfrac{6}{20}$\\
		$V_1 B_2$ & \vspace{5pt}$\dfrac{3}{5}\cdot\dfrac{2}{4}=\dfrac{6}{20}$\\
		$V_1 V_2$ & \vspace{5pt}$\dfrac{3}{5}\cdot\dfrac{2}{4}=\dfrac{6}{20}$\\ \midrule
		Total & 1 \\ \bottomrule
	\end{tabular}
\end{sidepicture}


\begin{example}
	Uma urna contém duas bolas brancas e três vermelhas. Sorteamos ao acaso duas delas, sem reposição. Assim, tem-se os possíveis resultados na Figura \ref{picture:prob_independ_ex1} e na Tabela \ref{table:prob_independ_ex1}. Deste modo,
	\[
		P(B_2\setminus V_1)=\frac{2}{4}\text{,}
	\]
	e
	\[
		P(B_2\cap V_1)=P(B_2\setminus V_1)P(V_1)=\frac{2}{4}\cdot \frac{3}{5}=\frac{6}{20}\text{,}
	\]
	onde, nesse caso, os eventos são dependentes.
\end{example}

\newpage

\begin{sidepicture}{0cm}{picture}{AAAA}
	\label{picture:prob_independ_ex2}
	\begin{tikzpicture}[
		grow=right,
		level 1/.style ={level distance=2cm, sibling distance=2.6cm},
		level 2/.style ={level distance=2cm, sibling distance=1.2cm},
	]
		\node [prob_tree_origin] {}
		child {
			node [prob_tree_node] {$V_1$}
			child {
				node [prob_tree_node] {$V_2$}
				edge from parent
				node[below, prob_tree_text]  {$\frac{3}{5}$}
			} child {
				node [prob_tree_node] {$B_2$}
				edge from parent
				node[above, prob_tree_text]  {$\frac{2}{5}$}
			}
			edge from parent
			node[below, prob_tree_text]  {$\frac{3}{5}$}
		} child {
			node [prob_tree_node] {$B_1$}
			child {
				node [prob_tree_node]{$V_2$}
				edge from parent
				node[below, prob_tree_text]  {$\frac{3}{5}$}
			} child {
				node [prob_tree_node]{$B_2$}
				edge from parent
				node[above, prob_tree_text]  {$\frac{2}{5}$}
			}
			edge from parent
			node[above, prob_tree_text]  {$\frac{2}{5}$}
		};
	\end{tikzpicture}
\end{sidepicture}

\begin{sidepicture}{-6cm}{table}{AAAA}
	\label{table:prob_independ_ex2}
	\begin{tabular}{cc}\\\toprule
		Resultados & Probabilidades\\  \midrule
		$B_1 B_2$ & \vspace{5pt}$\dfrac{2}{5}\cdot\dfrac{2}{5}=\dfrac{4}{25}$ \\ 
		$B_1 V_2$ & \vspace{5pt}$\dfrac{2}{5}\cdot\dfrac{3}{5}=\dfrac{6}{25}$\\
		$V_1 B_2$ & \vspace{5pt}$\dfrac{3}{5}\cdot\dfrac{2}{5}=\dfrac{6}{25}$\\
		$V_1 V_2$ & \vspace{5pt}$\dfrac{3}{5}\cdot\dfrac{3}{5}=\dfrac{9}{25}$\\ \midrule
		Total & 1 \\ \bottomrule
	\end{tabular}
\end{sidepicture}

\begin{example}
	Uma urna contém duas bolas brancas e três vermelhas. Sorteamos ao acaso duas delas, com reposição. Assim, tem-se os possíveis resultados na Figura \ref{picture:prob_independ_ex2} e na Tabela \ref{table:prob_independ_ex2}. Deste modo,
	\[
		P(B_2\setminus V_1)=\frac{2}{5}\text{,}
	\]
	e
	\[
		P(B_2\cap V_1)=P(B_2)P(V_1)=\frac{2}{5}\cdot \frac{3}{5}=\frac{6}{25}\text{,}
	\]
	onde, nesse caso, os eventos são independentes.
\end{example}

\section{Partição do Espaço Amostral}

Os eventos $E_1$, $E_2$, $\dots$, $E_k$ formam uma partição do espaço amostral, $S$, quando eles são mutuamente exclusivos e sua reunião é igual ao espaço amostral, isto é
\[
	E_1\cap E_2\cap \dots \cap E_k=\emptyset
\]
e
\[
	E_1\cup E_2\cup \dots \cup E_k=S\text{.}
\]

Um exemplo gráfico, com $k=4$, pode ser visto na Figura \ref{picture:prob_part_espaco_amostral}, onde $E_1\cap E_2 \cap E_3 \cap E_4=\emptyset$ e $E_1\cup E_2\cup E_3 \cup E_4=S$.

\begin{sidepicture}{0}{picture}{AAAA}
	\label{picture:prob_part_espaco_amostral}
	\begin{tikzpicture}
		\begin{scope}[local bounding box=graph1]
			\node[inner sep=0pt](large) at (0,0){
				\includegraphics[scale=0.35]{probabilidade/particao4.png}
			};
		\end{scope}
		\node[align=left, color=ocre] at ($(graph1.west)+(0.6cm, 1.7cm)$) {$E_1$};
		\node[align=left, color=ocre] at ($(graph1.west)+(0.8cm, -0.7cm)$) {$E_2$};
		\node[align=left, color=ocre] at ($(graph1.west)+(3.2cm, 1cm)$) {$E_3$};
		\node[align=left, color=ocre] at ($(graph1.west)+(3.2cm, -1.1cm)$) {$E_4$};
	\end{tikzpicture}
\end{sidepicture}

\section{Regra da Probabilidade Total}

A regra da probabilidade total é enunciado no Teorema \ref{theorem:prob_total}.

\begin{theorem}
\label{theorem:prob_total}
	Sejam $E_1$, $E_2$, $\dots$, $E_k$ uma partição de $S$ e $B$ um evento qualquer, então tem-se
	\[
		P(B) = P(B\setminus E_1)P(E_1)
		+ \cdots
		+ P(B\setminus E_k)P(E_k)
		\text{.}
	\]

	\begin{proof}
	O evento $B$ pode ser escrito como
	\[
		B = (B\cap E_1)\cup \dots \cup (B\cap E_k)\text{,}
	\]
	logo,
	\[
		P(B)=P(B\cap E_1)+\cdots+P(B\cap E_k)\text{,}
	\]
	onde, utilizando a regra do produto, tem-se
	\[
		P(B)=P(B\setminus E_1)P(E_1)+\cdots+P(B\setminus E_k)P(E_k)\text{.}
	\]
	\end{proof}
\end{theorem}

\section{Teorema de Bayes}

O Teorema de Bayes é enunciado no Teorema \ref{theorem:bayes}.

\begin{theorem}
	\label{theorem:bayes}
	Sejam $E_1$, $E_2$, $\dots$, $E_k$ uma partição do espaço amostral $S$ e, seja $B$ um evento qualquer de $S$ e, ainda, considere conhecidas $P(E_i)$ e $P(B\setminus E_i)$, para $i=1$, $2$, $\dots$, $k$, então a probabilidade de ocorrência de um dos eventos $E_i$, dado que ocorreu o evento $B$ é dada por
	\[
		P(E_i\setminus B)=\frac{
			P(B\setminus E_i)P(E_i)
		} {
			\displaystyle \sum_{j=1}^{k} P(B\setminus E_j)P(E_j)
		}
		\text{,}
	\]
	ou, ainda,
	\[
		P(E_i\setminus B)=\frac{
			P(B\setminus E_i)P(E_i)
		}{
			P(B\setminus E_1)P(E_1)
			+
			\cdots
			+
			P(B\setminus E_k)P(E_k)
		}
		\text{.}
	\]
	
	\begin{proof}
		Pela probabilidade condicional, tem-se que
		\[
			P(E_i\setminus B)=\frac{P(E_1\cap B)}{P(B)}\text{,}
		\]
		e, pela regra do produto, tem-se
		\[
			P(E_i\cap B) = P(B\cap E_i)=P(B\setminus E_i)P(E_i)\text{.}
		\]
		
		Pela regra da probabilidade total (Teorema \ref{theorem:prob_total}), tem-se
		\[
			P(B) = P(B\setminus E_1)P(E_1)+\cdots + P(B\setminus E_k)P(E_k)\text{,}
		\]
		portanto, tem-se
		\[
			P(E_i\setminus B)=\frac{
				P(B\setminus E_i)P(E_i)
			}{
				P(B\setminus E_1)P(E_1)
				+
				\cdots
				+
				P(B\setminus E_k)P(E_k)
			}
			\text{.}
		\]
	\end{proof}
\end{theorem}

\begin{pageWidthArea}
	\begin{example}
			Em uma fábrica de parafusos, as máquinas A, B e C produzem 25\%, 35\% e 40\% do total produzido, respectivamente. Da produção de cada máquina, 5\%, 4\% e 2\%, respectivamente, são parafusos defeituosos. Determine:\\
	
			\begin{enumerate}[label=\alph*), itemsep=10pt]
				\item A probabilidade de que o parafuso seja fabricado na máquina A, sabendo que ele é defeituoso.\\
				
				Representando por A, B e C as três máquinas e, ainda, por D os parafusos defeituosos. As probabilidades são $P(A)=0,25$, $P(B)=0,35$ e $P(C)=0,4$, $P(D\setminus A)=0,05$, $P(D\setminus B)=0,04$, $P(D\setminus C)=0,02$. Assim, pelo teorema de Bayes (Teorema \ref{theorem:bayes}), tem-se
				\begin{align*}
					P(A\setminus D) &=
					\frac{
						P(D\setminus A)P(A)
					}{
						P(D\setminus A)P(A)+P(D\setminus B)P(B)+P(D\setminus C)P(C)
					}\\
					&=\frac{0,05\cdot 0,25}{
						0,05\cdot 0,25
						+
						0,04\cdot 0,35
						+
						0,02\cdot 0,40
					}\\
					&=0,3623\text{.}
				\end{align*}
				
				\item A probabilidade de que o parafuso não seja defeituoso.\\
				
				Considerando as probabilidades $P(D^C\setminus A)=0,95$, $P(D^C\setminus B)=0,96$ e $P(D^C\setminus C)=0,98$. Deste modo,
				\begin{align*}
					P(D^C) &= P(D^C\setminus A)P(A)+P(D^C\setminus B)P(B)+P(D^C\setminus C)P(C)\\
					&=0,95\cdot 0,25+0,96\cdot 0,35+0,98\cdot 0,40\\
					&=0,9655\text{.}
				\end{align*}
				
				\item A probabilidade de que o parafuso seja fabricado na máquina C sabendo que ele não é defeituoso.\\
				
				Fazendo
				\begin{align*}
					P(C\setminus D^C)&=\frac{P(D^C\setminus C)P(C)}{P(D^C)}\\
					&=\frac{0,98\cdot 0,40}{0,9655}\\
					&=0,4060\text{.}
				\end{align*}
			\end{enumerate}
	\end{example}
\end{pageWidthArea}